\section{Arquitetura}
\label{sec:arquitetura}

A Figura~\ref{fig:arquitetura} mostra o fluxo de dados. Os sensores são lidos pelo ESP32, que detecta eventos localmente e os envia por HTTP à API. A API valida, classifica e grava as leituras no MongoDB; a aplicação web consulta a API periodicamente. Os caminhos tracejados são usados apenas na configuração: a descoberta do servidor por UDP e o pareamento Wi-Fi pelo ponto de acesso da placa.

\begin{figure*}[!t]
\centering
\begin{tikzpicture}[
  font=\footnotesize,
  >={Stealth[length=2.2mm]},
  bloco/.style={draw, rounded corners=2pt, thick, align=center, fill=white, inner sep=3pt},
  sensor/.style={bloco, fill=orange!12, minimum width=2.6cm, minimum height=0.9cm},
  esp/.style={bloco, fill=blue!8, minimum width=3.0cm, minimum height=3.7cm},
  srv/.style={bloco, fill=green!10, minimum width=3.0cm, minimum height=1.6cm},
  db/.style={bloco, fill=gray!15, minimum width=2.0cm, minimum height=1.6cm},
  rot/.style={font=\scriptsize, align=center, inner sep=1.5pt}
]
  % Sensores
  \node[sensor] (mic)   at (0, 1.3)  {Microfone\\eletreto + LM393};
  \node[sensor] (sonar) at (0, 0)    {Sonar\\LV-MaxSonar-EZ1};
  \node[sensor] (cap)   at (0, -1.3) {Capacitivo\\CR18-8DP};

  % ESP32
  \node[esp] (esp) at (5.0, 0) {\textbf{ESP32 (ESP-IDF)}\\[3pt]
     \textit{task} do microfone\\(núcleo 1)\\[2pt]
     laço principal:\\Wi-Fi, sonar, botão\\[2pt]
     fila + \textit{api\_worker}\\[2pt]
     NVS (configuração)};

  % API e banco
  \node[srv] (api) at (10.3, 0) {\textbf{API}\\Express 5 + Zod\\classificador de sons\\indicadores diários};
  \node[db]  (db)  at (14.4, 0) {\textbf{MongoDB}\\leituras\\dispositivos\\placas};

  % Web
  \node[srv, fill=purple!8] (web) at (10.3, -3.3) {\textbf{Aplicação web}\\React + TanStack Query\\planta, painel, saúde};

  % Sensores -> ESP32
  \draw[->, thick] (mic.east)   -- node[rot, above] {ADC1 + GPIO} (mic.east -| esp.west);
  \draw[->, thick] (sonar.east) -- node[rot, above] {UART2} (sonar.east -| esp.west);
  \draw[->, thick] (cap.east)   -- node[rot, above] {GPIO25} (cap.east -| esp.west);

  % ESP32 -> API
  \draw[->, thick] ([yshift=6mm]esp.east) -- node[rot, above] {HTTP POST\\(leituras e\\\textit{heartbeat})} ([yshift=6mm]esp.east -| api.west);
  \draw[<->, dashed] ([yshift=-6mm]esp.east) -- node[rot, below] {UDP 3333\\descoberta} ([yshift=-6mm]esp.east -| api.west);

  % API <-> MongoDB
  \draw[<->, thick] (api.east) -- node[rot, above] {Mongoose} (db.west);

  % Web <-> API
  \draw[<->, thick] (web.north) -- node[rot, right] {HTTP (REST)\\a cada 3\,s} (api.south);

  % Web -> ESP32 (pareamento)
  \draw[->, dashed] (web.west) -| node[rot, pos=0.3, below] {pareamento Wi-Fi via\\\texttt{PetCare-Setup} (10.73.73.1)} (esp.south);
\end{tikzpicture}
\caption{Arquitetura do PetCare. Setas contínuas: fluxo de dados em operação; tracejadas: fluxos de configuração (descoberta e pareamento).}
\label{fig:arquitetura}
\end{figure*}

\subsection{Contrato entre o ESP32 e a API}

O ponto central da integração é a rota \texttt{POST /api/readings}. O corpo é validado com Zod e aceita uma leitura ou um lote de até 100 leituras (objeto com o campo \texttt{readings} contendo um vetor). Os campos são: \texttt{sensorKey} (obrigatório), \texttt{value} (obrigatório), \texttt{recordedAt}, \texttt{durationMs} (até 60\,s), \texttt{kind}, \texttt{features} e \texttt{envelope} (até 128 valores entre 0 e 150). Exemplos de corpo enviados pelo firmware:

\begin{lstlisting}
{ "sensorKey": "esp32-01:proximity", "value": 1 }

{ "sensorKey": "esp32-01:mic", "value": 48.3,
  "durationMs": 920,
  "features": { "f0Hz": 610, "f0StartHz": 560,
    "f0EndHz": 650, "voicedRatio": 0.72,
    "zcr": 0.081, "bursts": 1 },
  "envelope": [31, 44, 48, 47, 40, 33] }
\end{lstlisting}

Os valores do segundo exemplo são ilustrativos do formato, e não uma medição. Quando a API tem a variável \texttt{DEVICE\_TOKEN} definida, a placa deve enviar o mesmo valor no cabeçalho \texttt{x-device-token}; caso contrário, a API responde 401. A API também atualiza o horário da última comunicação da placa, usado para mostrá-la como \textit{online} ou \textit{offline}; um dispositivo é considerado \textit{online} se enviou leitura nos últimos 5 minutos.

\subsection{API}

A API é organizada em módulos por domínio (\texttt{devices}, \texttt{readings}, \texttt{pairing}, \texttt{pets}, \texttt{insights}), cada um com modelo, esquemas de validação, rotas e serviço. Além da ingestão, ela oferece o cadastro dos itens da planta (nome, tipo, \textit{sensorKey}, posição em porcentagem da planta e capacidade do pote), o histórico por dispositivo, o estado calculado de cada item (último nível, eventos no dia, \textit{online}), a lista de placas pareadas e um servidor UDP que responde à descoberta.

\paragraph{Classificação das vocalizações.}
Quando a leitura do microfone não traz o tipo (\texttt{kind}), a API o estima por regras aplicadas em ordem (Tabela~\ref{tab:classificador}), inspiradas na literatura de acústica felina, que relaciona duração, $f_0$ e contorno de entonação dos miados ao contexto \citep{schotz2019}. Manter as regras na API permite ajustar os limiares sem regravar o firmware. \pendente{Os limiares não foram validados com gravações rotuladas; ver Seção~\ref{sec:resultados}.}

\begin{table}[!ht]
\caption{Regras do classificador de vocalizações, aplicadas em ordem (\texttt{vocalization.classifier.ts}).}
\label{tab:classificador}
\centering
\small
\begin{tblr}{
  colspec={X[0.9,l] X[2,l]},
  row{1}={font=\bfseries},
  hline{1,Z}={0.8pt}, hline{2}={0.4pt},
  rowsep=1pt
}
Tipo & Regra \\
Sibilo/rosnado & fração tonal $< 0{,}3$ e ZCR $\geq 0{,}25$; ou $f_0 < 250$\,Hz e duração $> 0{,}8$\,s \\
Tagarelar & $\geq 3$ rajadas com duração média por rajada $< 200$\,ms \\
Trinado & duração $< 0{,}6$\,s e $f_0$ subindo mais de 300\,Hz/s \\
Uivo & duração $\geq 1{,}5$\,s e (volume $\geq 50$ ou variação de $f_0 \geq 250$\,Hz); sem tom disponível, só a duração \\
Miado longo & duração $\geq 0{,}8$\,s \\
Miado curto & demais casos \\
\end{tblr}
\end{table}

\paragraph{Indicadores diários.}
O módulo \texttt{insights} agrega as leituras por dia: consumo de água e ração (quedas de nível de pelo menos 2 pontos percentuais, após suavização em janelas de 5 minutos, com subidas de 5 pontos ou mais tratadas como reabastecimento), número e duração das visitas à caixa de areia (permanências acima de 30 minutos são descartadas como falha de detecção) e vocalizações por tipo, inclusive noturnas (22h--6h). A aba de saúde compara esses valores com faixas de referência e com a própria linha de base do gato. Como o CR18-8DP só fornece 0\,\% ou 100\,\%, o cálculo de consumo, pensado para nível contínuo, fica restrito a transições cheio/vazio (Seção~\ref{sec:resultados}).

\subsection{Aplicação web}

A aplicação tem cinco páginas: \textbf{planta baixa} (itens arrastáveis sobre a planta da casa, com estado em tempo quase real), \textbf{painel} (gráficos por dispositivo, incluindo a forma de onda do envelope de cada som colorida pelo tipo detectado), \textbf{saúde} (consumo, caixa de areia e vocalizações, com perfil do gato), \textbf{dispositivos} (tabela de cadastro) e \textbf{configurações} (pareamento guiado da placa e lista de placas pareadas). Os dados são atualizados por consulta periódica a cada 3\,s. \pendente{Inserir capturas de tela da planta baixa, do painel e da aba de saúde.}
